\documentclass[12pt,a4paper,UTF8]{ctexart}
\usepackage[a4paper,top=22mm,bottom=22mm,left=22mm,right=22mm,headheight=15pt]{geometry}
\usepackage{amsmath,amssymb,mathtools,bm}
\usepackage{graphicx,xcolor,booktabs,longtable,array,enumitem,fancyhdr,caption,float,microtype,setspace}
\usepackage[hidelinks,unicode]{hyperref}
\usepackage[most]{tcolorbox}
\definecolor{ContestBlue}{HTML}{17365D}
\hypersetup{pdftitle={第084套},pdfauthor={}}
\setstretch{1.24}
\setlength{\parindent}{2em}
\setlength{\parskip}{0.35em}
\setlength{\emergencystretch}{2em}
\setlist{itemsep=0.35em,topsep=0.4em}
\captionsetup{font=small,labelfont=bf,justification=centering}
\setcounter{secnumdepth}{0}
\pagestyle{fancy}
\fancyhf{}
\fancyhead[L]{\small 第084套}
\fancyfoot[C]{\thepage}
\renewcommand{\headrulewidth}{0.4pt}
\ctexset{
  section={format=\Large\bfseries\color{ContestBlue},beforeskip=1.2em,afterskip=0.7em},
  subsection={format=\large\bfseries\color{ContestBlue},beforeskip=1.0em,afterskip=0.5em},
  subsubsection={format=\normalsize\bfseries,beforeskip=0.8em,afterskip=0.35em}
}
\makeatletter
\newcommand{\tightlist}{\setlength{\itemsep}{0pt}\setlength{\parskip}{0pt}}
\newcommand{\maxwidth}{\ifdim\Gin@nat@width>\linewidth\linewidth\else\Gin@nat@width\fi}
\newcommand{\maxheight}{\ifdim\Gin@nat@height>0.72\textheight 0.72\textheight\else\Gin@nat@height\fi}
\newcommand{\pandocbounded}[1]{\begingroup\setkeys{Gin}{width=\maxwidth,height=\maxheight,keepaspectratio}#1\endgroup}
\makeatother
\begin{document}
\begin{titlepage}
  \centering
  \vspace*{58mm}
  {\Huge\bfseries\color{ContestBlue} 第084套\par}
  \vfill
  \vspace*{28mm}
\end{titlepage}
\section{第 1 题}\label{ux7b2c-1-ux9898}

如图，一只轻质桶内装有密度为 \(\rho_0\) 的液体。液体中有一密度为
\(\rho\)、体积为 \(V\)
的物块，物块通过理想滑轮和轻绳与桶相连。静止时物块部分浸入液体，既未完全浸没也未离开液面，并且与桶底、桶壁均不接触。已知浸入液体中的物块体积为

\[
V_1=\alpha V.
\]

\begin{enumerate}
\def\labelenumi{\arabic{enumi}.}
\tightlist
\item
  求桶内液体体积与物块体积之比 \(V_0/V\)；
\item
  在竖直方向给系统一个小振动。忽略液体流动的动能，桶的横截面积很大；物块为边长
  \(l\) 的立方体。求系统的小振动周期。
\end{enumerate}

\begin{figure}[H]
\centering
\includegraphics[width=0.548\linewidth,height=\textheight,keepaspectratio,alt={图 1-1　液体桶、固定滑轮、动滑轮与部分浸入物块的平衡状态。}]{../figures/question-01.pdf}
\caption*{图 1-1　液体桶、固定滑轮、动滑轮与部分浸入物块的平衡状态。}
\end{figure}

\newpage

\section{第 2 题}\label{ux7b2c-2-ux9898}

正八面体的每条棱均接入一个阻值为 \(R\)
的电阻，构成如图所示的六节点电阻网络。

\begin{enumerate}
\def\labelenumi{\arabic{enumi}.}
\tightlist
\item
  求顶点 \(A\)、\(B\) 之间的等效电阻；
\item
  求相邻顶点 \(A\)、\(C\) 之间的等效电阻；
\item
  从 \(A\) 点注入电流 \(I\)，从 \(C\) 点注入电流 \(2I\)，并从 \(B\)
  点流出电流 \(3I\)。求电势差 \(U_{AD}\)。
\end{enumerate}

\begin{figure}[H]
\centering
\includegraphics[width=0.9\linewidth,height=\textheight,keepaspectratio,alt={图 2-1　正八面体十二条棱组成的空间电阻网络。}]{../figures/question-02.pdf}
\caption*{图 2-1　正八面体十二条棱组成的空间电阻网络。}
\end{figure}

\newpage

\section{第 3 题}\label{ux7b2c-3-ux9898}

本题只考虑纸面内的二维近轴成像，并假定即使入射角较大，球面折射仍可使用近轴球面折射公式。

如图，\(OA\)、\(OB\) 为两面平面镜，\(OA=OB=l\)。圆弧 \(ACB\) 的半径为

\[
R=\alpha l,\qquad \alpha\gg1,
\]

扇形区域 \(AOB\) 内充满折射率为 \(n\) 的介质。光源 \(S\) 位于
\(\angle AOB\) 的角平分线上，且

\[
OS=\sqrt2\,l.
\]

计算时保留 \(1/\alpha\) 的一阶项。

\begin{enumerate}
\def\labelenumi{\arabic{enumi}.}
\tightlist
\item
  当 \(OA\perp OB\) 时，求光线经圆弧折射、\(OA\)、\(OB\)
  反射并再次折射后的像；
\item
  当 \(\angle AOB=120^\circ\) 时，求光线经折射、\(OA\)
  反射并再次折射后的像；
\item
  当 \(\angle AOB=120^\circ\) 时，求经 \(OA\)、\(OB\)
  两次反射再折射后形成的像。
\end{enumerate}

\begin{figure}[H]
\centering
\includegraphics[width=0.56\linewidth,height=\textheight,keepaspectratio,alt={图 3-1　圆弧折射面、双平面镜与角平分线上的光源。}]{../figures/question-03.pdf}
\caption*{图 3-1　圆弧折射面、双平面镜与角平分线上的光源。}
\end{figure}

\newpage

\section{第 4 题}\label{ux7b2c-4-ux9898}

一根轻绳的一端固定在质量为 \(M=\alpha m\)、半径为 \(R\)
的匀质圆盘边缘，并在圆盘上缠绕若干圈。轻绳的另一端绕过固定在墙上的圆柱后，与质量为
\(m\) 的物体相连。开始时物体 \(m\)
与圆盘的最低点等高，它们各自正下方距离 \(H\)
处分别有一位老师午睡。现将系统由静止释放。除第（6）问外，绳与固定圆柱之间光滑无摩擦，且圆盘上的绳足够长。

\begin{enumerate}
\def\labelenumi{\arabic{enumi}.}
\tightlist
\item
  在任何物体碰到下方老师以前，求绳中的张力；
\item
  圆盘质心是否可能向上运动；
\item
  要使左侧物体 \(m\) 向上运动，\(\alpha\) 应满足什么条件；
\item
  已知 \(\alpha\) 的取值能使两位老师最终都被碰到，判断哪一位先被碰到；
\item
  承接第（4）问，若认为圆盘一侧先发生的碰撞不影响物体 \(m\)
  的后续运动，求左侧老师经过多长时间被碰到；
\item
  若绳与固定圆柱之间的动摩擦因数为 \(\mu\)，求使物体 \(m\)
  向上运动所需的 \(\alpha\) 的范围。
\end{enumerate}

\begin{figure}[H]
\centering
\includegraphics[width=0.562\linewidth,height=\textheight,keepaspectratio,alt={图 4-1　悬挂物体、定滑轮与绕绳圆盘组成的动力学系统。}]{../figures/question-04.pdf}
\caption*{图 4-1　悬挂物体、定滑轮与绕绳圆盘组成的动力学系统。}
\end{figure}

\newpage

\section{第 5 题}\label{ux7b2c-5-ux9898}

\begin{enumerate}
\def\labelenumi{\arabic{enumi}.}
\tightlist
\item
  有一条均匀带电无限长直线，线电荷密度为 \(\lambda\)，求距直线为 \(r\)
  处的电场 \(\boldsymbol E_1(r)\)；
\item
  半径为 \(R\) 的均匀带电圆环位于 \(xOy\)
  平面内，圆心在原点，线电荷密度为
  \(\lambda\)。空间中无其他带电体。求轴线上 \((0,0,z)\) 处的电场
  \(\boldsymbol E_2(z)\)，并在 \(|z|\ll R\) 时保留关于 \(z/R\)
  的最低非零阶；
\item
  半径为 \(r\)、线电荷密度为 \(\lambda\) 的均匀带电圆环位于 \(xOy\)
  平面内，圆心在原点。另有一条与 \(z\) 轴平行、通过 \((a,0,0)\)
  的无限长直线，线电荷密度也为 \(\lambda\)。在 \(a>0\) 且 \(a\ne r\)
  时，求直线对圆环的总静电力。可以使用围道积分，也可以利用系统沿 \(z\)
  方向的平移对称性化简。
\end{enumerate}

\begin{figure}[H]
\centering
\includegraphics[width=0.9\linewidth,height=\textheight,keepaspectratio,alt={图 5-1　带电圆环与平行于其轴线的无限长直线电荷。}]{../figures/question-05.pdf}
\caption*{图 5-1　带电圆环与平行于其轴线的无限长直线电荷。}
\end{figure}

\newpage

\section{第 6 题}\label{ux7b2c-6-ux9898}

\begin{enumerate}
\def\labelenumi{\arabic{enumi}.}
\item
  一电子在相互正交的匀强电场 \(\boldsymbol E\) 和匀强磁场
  \(\boldsymbol B\)
  中运动。稳定后，电子所受洛伦兹力的时间平均值为零。求漂移速度的时间平均值
  \(\overline{\boldsymbol V}_{D1}\)，并约定
  \(\overline{\boldsymbol V}_{D1}\perp\boldsymbol B\)；
\item
  空间中只有大小为 \(B\) 的磁场，磁感线弯曲且曲率半径为 \(R\)。电荷量为
  \(-e\)、质量为 \(m\) 的电子在其中运动，沿磁场方向的速度为
  \(V_{\parallel}\)，垂直磁场方向的速度为 \(V_{\perp}\)，且
  \(V_{\parallel},V_{\perp}\ll eBR/m\)，相应回旋尺度远小于
  \(R\)。求磁场曲率引起的漂移速度；
\item
  已知在非均匀磁场

  \[
  \boldsymbol B=B_0(1+\alpha x)\hat{\boldsymbol z}
  \]

  中，以半径 \(r_L\) 做圆周运动的电子具有梯度漂移

  \[
  \overline{\boldsymbol V}_{D2}
  =-\frac{\alpha eB_0r_L^2}{2m}\hat{\boldsymbol y},
  \qquad r_L\ll\frac1{|\alpha|}.
  \]

  若弯曲磁场的大小随到曲率中心的距离 \(\rho\) 按

  \[
  B(\rho)=\frac{BR}{\rho}
  \]

  变化，求包括曲率漂移和梯度漂移在内的总漂移速度；
\item
  电子处于温度为 \(T\) 的热平衡状态，求总漂移速度的热平均值。取
\end{enumerate}

\[
m=9.11\times10^{-31}\,\mathrm{kg},\quad
e=1.60\times10^{-19}\,\mathrm C,
\]

\[
B=0.10\,\mathrm T,\quad
R=1.0\,\mathrm m,\quad
T=\frac{1\,\mathrm{eV}}{k_B},
\]

给出数值结果。

\begin{figure}[H]
\centering
\includegraphics[width=0.44\linewidth,height=\textheight,keepaspectratio,alt={图 6-1　弯曲磁场中的电子回旋运动与导引中心漂移。}]{../figures/question-06.pdf}
\caption*{图 6-1　弯曲磁场中的电子回旋运动与导引中心漂移。}
\end{figure}

\newpage

\section{第 7 题}\label{ux7b2c-7-ux9898}

如图，连通器的 \(A\)、\(B\) 两个上部容器横截面积分别为 \(4S\) 和
\(S\)，底部水平容器的横截面积为 \(S\)。连通器中装有密度为 \(\rho_1\)
的液体甲。初始时，甲液面到 \(A\) 容器口的距离为
\(H\)，且两侧液面到底部的距离足够大。装置处于恒温环境，外界大气压为

\[
p_0=\rho_1gh_0,
\qquad h_0<H.
\]

已知容器体积变化所引起的气体压强变化必须计入。

\begin{enumerate}
\def\labelenumi{\arabic{enumi}.}
\tightlist
\item
  从 \(B\) 端缓慢注入与甲不相溶、密度为 \(\rho_2=\rho_1/2\)
  的液体乙，乙的总体积为 \(SH\)，液体均不溢出。求稳定后 \(A\)、\(B\)
  两边上表面的高度差 \(h_1\)；
\item
  若先将 \(A\) 容器口封闭，再从 \(B\)
  端完成同样的注液，求稳定后两边上表面的高度差 \(h_2\)；
\item
  若初始时 \(A\)、\(B\) 两端甲液体的深度均为
  \(H/4\)，底部水平容器足够长，重复第（1）问，求高度差 \(h_3\)；
\item
  在第（3）问的浅液初态下，先封闭 \(A\)
  容器口，再重复第（2）问，求高度差 \(h_4\)。
\end{enumerate}

取

\[
h_0=2.82\,\mathrm m,
\qquad
H=3.51\,\mathrm m.
\]

\begin{figure}[H]
\centering
\includegraphics[width=0.723\linewidth,height=\textheight,keepaspectratio,alt={图 7-1　上部截面积不同、底部相连的连通器。}]{../figures/question-07.pdf}
\caption*{图 7-1　上部截面积不同、底部相连的连通器。}
\end{figure}

\newpage

\section{第 8 题}\label{ux7b2c-8-ux9898}

一种理想化轮子输送机的轨道由直线段 \(AB\)、半圆弧段 \(BC\) 和直线段
\(CD\) 组成，三段均位于纸面内，半圆弧半径为 \(r\)。沿轨道 \(ABCD\)
密排着 \(N\)
个轻轮子，每个轮子的轮轴固定，轮子可以绕自身轮轴无摩擦转动。重力方向垂直纸面向里。

在 \(A\) 处给一个质量为 \(m\)、绕质心转动惯量为 \(I\) 的小物体沿 \(AB\)
指向 \(B\) 的初速度
\(v_0\)。小物体与轮子之间的滑动摩擦因数很大；小物体尺寸很小，任一时刻至多与两个轮子接触。最终小物体沿
\(ABCD\) 从 \(A\) 到达 \(D\)。

\begin{enumerate}
\def\labelenumi{\arabic{enumi}.}
\tightlist
\item
  判断小物体到达 \(D\) 点时速度的方向；
\item
  求小物体到达 \(D\) 点时的速度大小。
\end{enumerate}

\begin{figure}[H]
\centering
\includegraphics[width=0.664\linewidth,height=\textheight,keepaspectratio,alt={图 8-1　直线---半圆---直线组成的轮子输送轨道。}]{../figures/question-08.pdf}
\caption*{图 8-1　直线---半圆---直线组成的轮子输送轨道。}
\end{figure}
\end{document}
